\chapter[Introducción]{Introducción}
\label{cp:introduction}

La idea de clasificar el suelo según lo que hacemos en él, ya sea para un cultivo, una zona residencial, un bosque o un polígono industrial, 
es algo que sirve para la planificación de diferentes tareas para la gestión de ciudades. Saber qué hay en cada metro cuadrado del territorio
es fundamental para decidir qué zonas naturales proteger, cómo planificar el crecimiento de las ciudades, cómo repartir las ayudas a los
agricultores o cómo actuar con rapidez ante una emergencia como incendios o inundaciones.  


Hasta ahora, en España, este trabajo de clasificación se revisaba manualmente cada vez que se buscaba actualizar el mapa de cobertura y usos de 
suelo, sobre una base de datos generada de manera automática apoyándose en múltiples fuentes. Así se han formado grandes bases de datos como
\acrfull{siose} a nivel nacional o \acrshort{corine} Land Cover a nivel europeo, con la parte de \acrshort{corine} Land Cover España disponible
junto a \acrshort{siose} en el centro de descargas del \acrfull{ign}. No obstante, esto trae consigo una dependencia del ojo humano, haciendo
que este proceso sea extremadamente lento, al revisar toda la expansión del territorio de esta forma, provocando que estos mapas se actualicen
completamente cada ciertos años, generando que cambios diarios en el uso del suelo, como expansiones de ciudades, de parques naturales, o cualquier
cambio del suelo, no se vean reflejados en los mapas hasta varios años más tarde.


Esta es una de las principales causas por las cuales surge este proyecto, junto con la ayuda de la \acrshort{ia} y satélites como la
constelación Sentinel-2 del programa europeo Copernicus, tenemos disponible un flujo constante de imágenes de alta resolución de la Tierra. 
Mediante la fusión de estos 2 recursos, podemos aspirar a automatizar en gran medida este proceso de clasificación, de manera similar a como 
lo hace BigEarthNet con otros países europeos. El plan y objetivo central de este trabajo es desarrollar un sistema semejante, pero diseñado
para clasificar el suelo español peninsular.


Este \acrfull{tfg} nace de una propuesta real y de una necesidad tecnológica planteada por la empresa \textbf{Xcalibur}, especializada en 
teledetección y mapeo avanzado del territorio. Este documento pretende explicar el flujo de trabajo realizado para la creación de los modelos,
partiendo desde la extracción de 6 macroparches con imágenes de Sentinel-2 extraídos de \acrfull{gee} hasta la evaluación de los resultados, pasando 
por todos los ajustes que son necesarios para cuadrar los datos gracias a la librería \textbf{Rasterio} con los mapas vectoriales ofrecidos por 
\acrshort{corine} Land Cover España bajo la proyección EPSG:25830 y la visualización de estos de manera masiva gracias a la herramienta 
\textbf{QGIS}. Para no trabajar con imágenes con miles de kilómetros cuadrados de área, se trocearon en archivos de extensión npy tanto para 
datos de entrada como de salida, bajo la nomenclatura de archivos \texttt{\_X.npy} (para las bandas 2,3,4 y 8 del satélite) y \texttt{\_Y.npy} 
(para las etiquetas de uso del suelo), organizados en matrices de $50\times 50$ píxeles.


Con el conjunto de datos preparado, implementamos, entrenamos y evaluamos dos modelos de redes neuronales basados en una estructura de \acrshort{CNN},
para la parte de obtención de las características de los parches. El primer modelo se trata de un modelo de clasificación en el cual un bloque 
convolucional alimenta a un \acrshort{mlp} con el fin de clasificar los parches entrantes asignándoles una probabilidad a cada una de las etiquetas y devolviendo la clase mayoritaria 
como resultado. El segundo modelo utiliza la parte convolucional como un encoder que genera un espacio latente, conectándolo con un bloque convolucional
inverso (el decoder) mediante una arquitectura \textbf{U-net} para devolver las etiquetas a nivel de píxel en vez de parche, lo cual permite predecir con mayor resolución los usos de suelo en 
cada parche. Para que las operaciones de la U-net sean exactas reduciremos los parches de $50\times 50$ a $48\times 48$ permitiendo un mejor 
flujo de los datos. Vimos que había un gran desbalanceo en España, habiendo muchas zonas naturales comparado con otras zonas, como
industriales o comerciales, se añadieron unos sesgos en la función de pérdida asociados a cada una de las clases, dando mayor importancia a las 
que aparecen menos. Finalmente, evaluamos los modelos graficando la evolución de métricas como la pérdida o la precisión, con la ayuda de \textbf{matplotlib} y matrices de confusión
para entender los resultados de nuestros modelos. 

\section{Contexto y Planteamiento del Problema}
\label{sec:contexto_problema}

Para comprender la necesidad de automatizar la cartografia del territorio, es necesario entender que el suelo es dinámico, por lo que cosas esenciales
en un país, como la ordenación del territorio y el diseño de políticas medioambientales, requieren una revisión sistemática y constante, ya que
fenómenos como la expansión urbana, la deforestación o los cambios drásticos en los cultivos, modifican el mapa periodicamente.


Bases de datos como \acrshort{siose} o \acrshort{corine} Land Cover han ayudado con la creación de mapas de usos y cobertura del suelo para la gestión
tanto a nivel español como europeo. Sin embargo, crear estos mapas conlleva un gran esfuerzo, requiriendo la revisión manual de los mapas generados 
automáticamente mediante fuentes disponibles como Catastro, \acrfull{sigpac} o la Información de la Declaración de Agricultores \cite{siose.es}.
Es por esto que los mapas a nivel nacional se actualicen cada 6 años, lo que no permite reflejar todos los cambios necesarios para un óptimo
manejo de los recursos, lo que perjudica a las empresas y entidades que tengan que trabajar con estas cartografías desactualizadas para tomar decisiones relevantes.


Por lo tanto, el problema que aborda este \acrshort{tfg} no es la falta de imágenes o fuentes para construir estos mapas, sino proponer una solución
al problema de la revisión completamente humana de estos, mediante la implementación de modelos de \acrshort{ia}, con el fin de actualizar estos
mapas con mayor frecuencia, buscando mantener la precisión.

\section{La oportunidad de la Inteligencia Artificial}
\label{sec:oportunidad_ia}

Para superar las limitaciones de los métodos manuales y adaptarse a la velocidad con la que cambia el entorno, la \acrshort{ia}, y más concretamente
el aprendizaje profundo (\textit{Deep Learning}), se presenta como una solución excelente. Esto es posible debido a, en primer lugar, la disponibilidad
de constelaciones de satélites, como es el caso de Sentinel-2, los cuales generan imágenes multiespectrales de la Tierra, y, en segundo lugar,
a la capacidad de las redes neuronales de analizar una cantidad masiva de datos, más concretamente las \acrshort{cnn}, por su diseño capaz de
analizar imágenes. 


Por un lado, Sentinel-2, aporta un flujo continuo de datos, con distintas resoluciones y bandas espectrales. Al estar actualizando los datos
con alta frecuencia, de todas las bandas que obtiene Sentinel-2 nosotros nos quedaremos con tan solo 4. Esta obtención actualizada y continua de
datos nos solventa el problema de la captura de datos, y dejar de depender de tantas fuentes como hasta ahora en \acrshort{siose}, ofreciendo un
conjunto de datos actualizado y fiable. Por otro lado, los algoritmos de aprendizaje profundo son capaces de procesar estas ingentes cantidades 
de datos y aprender a reconocer patrones en estos sobre el terreno, asignándole a cada píxel en los datos un uso de suelo específico.


Esta idea ya fue puesta realizada por parte de BigEarthNet, quienes, mediante parches de ambas constelaciones Sentinel (1 y 2) para unos
determinados países europeos consiguieron entrenar redes neuronales profundas con un nivel de precisión y automatización elevado. Siguiendo el
modelo internacional, este trabajo plantea el plan de trasladar esa misma filosofía al suelo español peninuslar, construyendo una herramienta
que facilite la clasificación de los usos de suelo de forma frecuente y automatizada.


\section{Motivación del proyecto}
\label{sec:motivacion}

Este \acrshort{tfg} surge ante una propuesta de la empresa \textbf{Xcalibur}, una empresa líder en el sector de la industria geofísica
aérea y mapeo de alta resolución del territorio \cite{xcalibur}, más concretamente de su departamento de \acrshort{ia} que buscaba entrenar un
modelo para la detección de los usos de suelo dadas las imágenes de la constelación Sentinel-2.


Dentro de su actividad, poder disponer de mapas actualizados sobre los usos de suelo y no tener que realizar vuelos con tanta frecuencia para 
poder obtener estos datos es de gran ayuda, al estar las principales bases de datos españolas, como \acrshort{corine} Land Cover o \acrshort{siose} desactualizadas.
Es por esto que plantearan desde Xcalibur construir un modelo que ayude en la construcción de mapas de usos de suelo actualizados con más frecuencia,
buscando obtener una precisión suficiente como para no dedicar tanto tiempo humano como se dedica actualmente en las revisiones manuales.    


\section{Objetivos del proyecto}
\label{sec:Objetivos}

Para dar una solución efectiva al problema planteado por Xcalibur, este trabajo se articula en torno a un objetivo principal, del cual surgen
unos más específicos.


\subsection{Objetivo General}

El objetivo de este proyecto es diseñar, implementar y evaluar un sistema de inteligencia artificial que permita automatizar, o al menos
facilitar, la creación de cartografías que clasifiquen y hagan una segmentación semántica de los usos del suelo de la España peninsular a partir
de imágenes multiespectrales de la constelación Sentinel-2, tomando como referencia metodológica BigEarthNet.

\subsection{Objetivos Específicos}
Para lograr la meta principal, el desarrollo del proyecto se divide en los siguientes objetivos específicos:

\begin{itemize}
    \item \textbf{Descarga de los datos geoespaciales:} Descargar del motor de \acrshort{gee} las 4 bandas de 10 metros de resolución de Sentinel-2 y descargar las coberturas oficiales de \acrshort{corine} Land Cover España.
    \item \textbf{Alineación cartográfica:} Sincronizar ambas fuentes de datos usando la librería \textit{Rasterio}, asegurando que las imágenes y las máscaras vectoriales con las etiquetas del uso del suelo coincidan bajo la proyección oficial EPSG:25830.
    \item \textbf{Generación física del dataset estructurado:} Fragmentar las 6 regiones descargadas de \acrshort{gee} en aproximadamente 1000000 de parches más manejables de $50\times50$ píxeles, exportándolos como archivos emparejados (\texttt{\_X.npy} y \texttt{\_Y.npy}), eliminando aquellos que contengan algún valor nulo y redefiniendo las etiquetas de usos de suelo, pasando de las 44 categorías de Land Cover a las 19 usadas en BigEarthNet.
    \item \textbf{Entrenamiento y evaluación de los modelos:} Implementar y comparar críticamente dos aproximaciones basadas en \acrshort{cnn}\begin{itemize}
        \item un clasificador a nivel de parche para predecir la clase dominante por parche.
        \item un modelo de segmentación semántica para predecir el uso de suelo a nivel de píxel, basado en la arquitectura U-Net.
    \end{itemize}. En las cuales reducimos el tamaño del parche a $48\times 48$ para facilitar la descomposición del parche dentro de la red.
    \item \textbf{Mitigación del desbalanceo de clases:} Diseñar e incorporar una función de pérdida avanzada basada en el Coeficiente de Dice Multiclase, aplicando sesgos y pesos específicos para corregir el desbalanceo de clases.
    \item \textbf{Análisis de los resultados:} Evaluar los resultados finales de ambos modelos, compararlos con BigEarthNet y proponer posibles mejoras a futuro.
\end{itemize}
